% Research development: exponential tightness of the first-passage mass
% up to a logarithmic exceptional set. To be read conditional on Lemmas
% \ref{lem:reduction}, \ref{lem:valuations}, \ref{lem:model} and the
% coupling in the proof of Proposition \ref{prop:clipped-calibration}.
\subsection*{Exponential tightness of the first-passage mass}
\label{subsec:tightness}
Lemma~\ref{lem:reduction} transfers not only means but also
distribution functions of first-passage observables, with an error that
does not depend on the observable. Combined with the coupling in the
proof of Proposition~\ref{prop:clipped-calibration}, this shows that the
moment condition just stated can fail only through an exceptional set of
logarithmic density $O((\log X)^{-\varepsilon_0})$. We keep the notation
of Proposition~\ref{prop:clipped-calibration} and
Corollary~\ref{cor:two-tails}, and put, for $\Lambda\ge0$ and $M\ge0$,
\[
 \begin{aligned}
 \mu_X(\Lambda)&=\lim_{N\to\infty}\frac1{\log N}
   \sum_{1\le w<N}\frac{\mathbf 1[\Phi_X(w^\flat)>\Lambda X]}{w},\\
 T_X(M)&=\lim_{N\to\infty}\frac1{\log N}
   \sum_{1\le w<N}\frac{(\Phi_X(w^\flat)-MX)_+}{w},
 \end{aligned}
\]
so that $T_X=T_X\bigl((\log X)^{\varepsilon_0/2}\bigr)$.

\begin{proposition}[Exponential tightness up to a logarithmic exceptional set]
\label{prop:tightness}
Assume the hypotheses of Proposition~\ref{prop:clipped-calibration}.
There are constants $C\ge1$ and $X_0$, independent of $X$ and $\Lambda$, such that
for all $X\ge X_0$ and all $\Lambda\ge0$ the limit $\mu_X(\Lambda)$
exists and
\begin{equation}\label{eq:tightness}
 \mu_X(\Lambda)\le C e^{-\Lambda/C}+C(\log X)^{-\varepsilon_0}.
\end{equation}
\end{proposition}

\begin{proof}
\emph{Step 1 (transfer).} Fix $X$ and $\Lambda$, and let
$\varphi=\mathbf 1[\Phi_X>\Lambda X]$. By Step~2 of the proof of
Theorem~\ref{thm:lower}, $\Phi_X(m)=\Phi_X(\operatorname{Pass}_x(m))$ for
all $x\ge x_0=X^{1+\beta}$ and all odd $m>x$ with $\tau_x(m)<\infty$ and
$\operatorname{Pass}_x(m)>\frac34x\exp(-(\log x)^{0.7})$, once $X$ is
large. Every function of $\Phi_X$ inherits this invariance. Hence
Lemma~\ref{lem:reduction} applies to $\varphi$: the limit
$\mu_X(\Lambda)=\delta_\varphi$ exists, and with $u=x_0^{\lambda_0}$,
$v=u^\alpha$ and $\mathbf N\sim\mathbf{Log}_{u,v}$,
\[
 \bigl|\mu_X(\Lambda)-\mathbb P(\Phi_X(\mathbf N)>\Lambda X)\bigr|
 \le C(\log u)^{-\varepsilon_0}\le C(\log X)^{-\varepsilon_0}.
\]
The constant of Lemma~\ref{lem:reduction} does not depend on
$\varphi$, hence not on $\Lambda$.

\emph{Step 2 (coupling).} Couple $\mathbf N$, $U$ and $(g_j)$ as in the
proof of Proposition~\ref{prop:clipped-calibration}, with the good event
$G$, $\mathbb P(G^c)\le CX^{-1/C}$, and the upper comparison
\eqref{eq:calib-upper-coupling}. For large $X$ we have
$1+C\varepsilon\le2$, so that
\[
 \mathbb P(\Phi_X(\mathbf N)>\Lambda X)
 \le CX^{-1/C}+\mathbb P(Z^+>\Lambda X/2).
\]

\emph{Step 3 (exponential moments of local counts).} For $q\le s$ let
$N_q$ be the number of pairs $(j,i)$ with $j\ge1$, $0\le i<g_j$ and
$z_{j,i}\in(q-1,q]$. We claim that there are absolute constants
$\theta_0>0$ and $K_0<\infty$ with
\begin{equation}\label{eq:local-exp-moment}
 \mathbb E\,e^{\theta_0N_q}\le K_0\qquad\text{for all }q\le s.
\end{equation}
As in the proof of Lemma~\ref{lem:model}(b), let $j_1$ be the first
$j\ge1$ with $U-D_j\le q$, and put $b=1+\log\frac32$; then
\[
 N_q\le g_{j_1}+Y,\qquad
 Y=\sum_{k\ge1}g_{j_1+k}\,\mathbf 1\bigl[D_{j_1+k-1}-D_{j_1}<b\bigr].
\]
Conditionally on $U$, $j_1$ is a stopping time for $(g_j)$. By the strong
Markov property, $(g_{j_1+k})_{k\ge1}$ is an i.i.d.\ sequence with the
law of $g_1$, independent of $U,g_1,\ldots,g_{j_1}$; thus $Y$ is
independent of $g_{j_1}$ and has the law of
$Y_0=\sum_{k\ge1}g_k\mathbf 1[D_{k-1}<b]$.

\emph{The overshooting valuation.} Let $t=U-q\ge0$. The index $j$ is the
first with $D_j\ge t$ and $g_j=v$ only if
$t-v\log2+\log3\le D_{j-1}<t$, an interval of length less than $v$, and
$g_j$ is independent of $D_{j-1}$. By the visit bound in
Lemma~\ref{lem:model}(ii),
\[
 \mathbb P(g_{j_1}=v\mid U)\le 2^{-v}\,
 \mathbb E\#\{j\ge1:\ t-v\log2+\log3\le D_{j-1}<t\}
 \le C(1+v)2^{-v},
\]
so $\mathbb E[e^{\theta g_{j_1}}\mid U]\le C\sum_{v\ge1}(1+v)(e^\theta/2)^v$,
which is bounded uniformly in $U$ for every $\theta<\log2$.

\emph{The returns after the crossing.} Let
$\kappa=1+\sup\{k\ge0:D_k<b\}$, which is finite almost surely. Then
$Y_0\le S_\kappa$ with $S_k=g_1+\dots+g_k$, and by Chernoff's bound in
Lemma~\ref{lem:model}(ii),
$\mathbb P(\kappa\ge k+1)\le\sum_{k'\ge k}e^{b/2}\varrho^{k'}
\le C\varrho^{k}$. Put
$\psi(\theta)=\mathbb E\,e^{\theta g_1}=e^\theta/(2-e^\theta)$. The
Cauchy--Schwarz inequality gives
\[
 \mathbb E\,e^{\theta Y_0}
 \le\sum_{k\ge1}\mathbb E\bigl[e^{\theta S_k}\mathbf 1[\kappa=k]\bigr]
 \le\sum_{k\ge1}\psi(2\theta)^{k/2}\,\mathbb P(\kappa\ge k)^{1/2}
 \le C\sum_{k\ge1}\bigl(\psi(2\theta)\varrho\bigr)^{k/2},
\]
which is finite as soon as $\psi(2\theta)\varrho<1$. Since $\psi(0)=1$
and $\varrho<1$, this holds for all small $\theta>0$. Choosing such a
$\theta_0<\log 2$ and combining the two bounds by independence proves
\eqref{eq:local-exp-moment}.

\emph{Step 4 (conclusion).} Every point counted in $Z^+$ satisfies
$z_{j,i}\le\ell_+\le s$, and $e^{\ell_+}=Xe^{C\varepsilon}\le2X$. Hence
\[
 \frac{Z^+}{X}\le2\sum_{i\ge0}e^{-i}N_{\ell_+-i}
 =\frac{2}{1-e^{-1}}\sum_{i\ge0}w_iN_{\ell_+-i},
 \qquad w_i=(1-e^{-1})e^{-i}.
\]
The weights $w_i$ add up to $1$. With
$\theta_1=\theta_0(1-e^{-1})/2$, convexity of the exponential and
\eqref{eq:local-exp-moment} give
\[
 \mathbb E\exp\bigl(\theta_1Z^+/X\bigr)
 \le\sum_{i\ge0}w_i\,\mathbb E\exp\bigl(\theta_0N_{\ell_+-i}\bigr)\le K_0,
\]
and Markov's inequality gives
$\mathbb P(Z^+>\Lambda X/2)\le K_0e^{-\theta_1\Lambda/2}$.
Together with Steps~1 and~2,
\[
 \mu_X(\Lambda)\le C(\log X)^{-\varepsilon_0}+CX^{-1/C}
   +K_0e^{-\theta_1\Lambda/2},
\]
which is \eqref{eq:tightness} after enlarging $C$.
\end{proof}

\noindent
In particular, this yields tightness in probability,
$\lim_{M\to\infty}\limsup_{X\to\infty}\mu_X(M)=0$.
It does \emph{not} yield uniform integrability of $\Phi_X/X$:
the exceptional probability in \eqref{eq:tightness} is independent
of $M$ and can carry a non-negligible weighted contribution.


\begin{corollary}[Where the clipping tail can live]
\label{cor:extreme-tail}
Assume the hypotheses of Proposition~\ref{prop:clipped-calibration}, and
let $C$ and $X_0$ be as in Proposition~\ref{prop:tightness}.
\begin{enumerate}[label=(\alph*)]
 \item For $X\ge X_0$ and $0\le L\le M$, the means $T_X(L)$ and $T_X(M)$
 exist, and
 \[
  T_X(M)\le T_X(L)\le T_X(M)+C^2Xe^{-L/C}+CXM(\log X)^{-\varepsilon_0}.
 \]
 \item If $M_X\ge(\log X)^{\varepsilon_0/2}$ and
 $M_X=o\bigl((\log X)^{\varepsilon_0}\bigr)$, then
 $T_X=T_X(M_X)+o(X)$. In particular, $T_X=o(X)$ if and only if
 $T_X(M_X)=o(X)$.
 \item Suppose that there are $p>1$, $A<(p-1)\varepsilon_0$ and $K$
 such that, for all large $X$,
 \[
  \limsup_{N\to\infty}\frac1{\log N}\sum_{1\le w<N}\frac1w
  \Bigl(\frac{\Phi_X(w^\flat)}X\Bigr)^p\le K(\log X)^A .
 \]
 Then $T_X=O\bigl(X(\log X)^{-\kappa}\bigr)$ with
 $\kappa=\bigl((p-1)\varepsilon_0-A\bigr)/p>0$.
\end{enumerate}
\end{corollary}

\begin{proof}
Let $B_X=\sum_{a\in W_X}a$. By Step~1 of the proof of
Theorem~\ref{thm:lower}, $0\le\Phi_X\le B_X$. Lemma~\ref{lem:reduction},
applied to the bounded invariant function $(\Phi_X-MX)_+/B_X$, gives the
existence of $T_X(M)$.

(a) The first inequality holds because $(t-MX)_+\le(t-LX)_+$. For
$t\ge0$,
\[
 (t-LX)_+-(t-MX)_+=X\int_L^M\mathbf 1[t>\lambda X]\,d\lambda .
\]
We insert $t=\Phi_X(w^\flat)$, multiply by $1/w$, sum over $w<N$ and
divide by $\log N$. For each $N$ the finite sum and the integral may be
interchanged. The resulting integrands converge to $\mu_X(\lambda)$ for
every $\lambda$ by Proposition~\ref{prop:tightness}, and they are bounded
by $2$ for large $N$. By dominated convergence and \eqref{eq:tightness},
\[
 T_X(L)-T_X(M)=X\int_L^M\mu_X(\lambda)\,d\lambda
 \le X\int_L^\infty Ce^{-\lambda/C}\,d\lambda+CX(M-L)(\log X)^{-\varepsilon_0}.
\]

(b) Apply (a) with $L=(\log X)^{\varepsilon_0/2}$, for which
$T_X(L)=T_X$, and $M=M_X$:
$0\le T_X-T_X(M_X)\le C^2Xe^{-(\log X)^{\varepsilon_0/2}/C}
+CXM_X(\log X)^{-\varepsilon_0}=o(X)$.

(c) For $t\ge0$ and $M>0$ we have $(t-MX)_+\le t^p(MX)^{1-p}$, hence
$T_X(M)\le KX(\log X)^AM^{1-p}$. Let $L=(\log X)^{\varepsilon_0/2}$. If
$(A+\varepsilon_0)/p\ge\varepsilon_0/2$, apply (a) with
$M=(\log X)^{(A+\varepsilon_0)/p}\ge L$; then
$KX(\log X)^AM^{1-p}$ and $CXM(\log X)^{-\varepsilon_0}$ are both
$O(X(\log X)^{-\kappa})$. Otherwise $A<\varepsilon_0(p/2-1)$, and the
bound with $M=L$ gives directly
$T_X\le KX(\log X)^{A-(p-1)\varepsilon_0/2}$, where
$A-(p-1)\varepsilon_0/2<-\kappa$. In both cases the term
$C^2Xe^{-L/C}$ is $O(X(\log X)^{-\kappa})$.
\end{proof}

\begin{remark}[What remains, and what the numerics can see]
\label{rem:tail-barrier}
Proposition~\ref{prop:tightness} says that actual orbits inherit the
exponential tail of the model except on an exceptional set $E_X$ of
logarithmic density $O((\log X)^{-\varepsilon_0})$. The exponent
$\varepsilon_0$ comes from \eqref{eq:ggm_hit} and \eqref{eq:ggm_tv}; it
is small and not explicit. On $E_X$ the only available bound is
$\Phi_X\le B_X<X^2$. Thus the trivial estimate of the contribution of
$E_X$ to $T_X$ is $X^2(\log X)^{-\varepsilon_0}$, which misses the
target $o(X)$ by a factor $X(\log X)^{-\varepsilon_0}$.
Corollary~\ref{cor:extreme-tail}(b) isolates exactly this: the clipping
tail vanishes if and only if the orbits with $\Phi_X>M_XX$ carry mass
$o(X)$, for one (equivalently, every) choice of $M_X$ between
$(\log X)^{\varepsilon_0/2}$ and $o((\log X)^{\varepsilon_0})$.
By definition, $\Phi_X(m)=0$ if the orbit of $m$ never reaches
$[1,X_1]$. This does not exclude contributions from orbits that
first descend below $X_1$ and later become unbounded or enter a
cycle above $X_1$. The obstruction for the clipping tail consists of
finite initial excursions that, before their first passage below $X_1$,
collect more than $X(\log X)^{\varepsilon_0/2}$ in $W_X$.
Excluding them requires information about individual orbits beyond the
$2$-adic horizon of Lemma~\ref{lem:valuations} and beyond the
total-variation scale of \eqref{eq:ggm_tv}; return certificates that
depend only on a residue modulo a fixed power of $2$ cannot provide it
(Proposition~\ref{prop:residue-obstruction}).

Numerically, the typical part behaves exactly as in the model (Run~I).
The program \path{tail_check.c} draws odd starts logarithmically from the
window $[X^{1+\beta},X^{(1+\beta)^2}]$ used above and computes
$\Phi_X/X$ along the actual orbit and along the random walk of
Lemma~\ref{lem:model}; starts whose orbits leave $128$-bit arithmetic
are recomputed exactly. For $c\in\{1,-1,5\}$,
$X\in\{10^{12},10^{20},10^{30}\}$ and $\beta\in\{0.05,0.005\}$, with
$4\cdot10^6$ starts per configuration ($2\cdot10^7$ for $c=1$,
$X=10^{30}$, $\beta=0.05$), the means and second moments of $\Phi_X/X$
agree with the model within the sampling error. For example, for $c=1$,
$X=10^{30}$ and $\beta=0.05$ the mean is $6.714$ against $6.712$ and the
second moment $69.33$ against $69.27$. The tail probabilities
$\mathbb P(\Phi_X/X>\Lambda)$ agree in all $18$ configurations with
$|z|\le2.1$. The sampled tails are consistent with exponential decay, with an
empirically fitted rate of about $0.21$ for $\beta=0.05$ and
$X\ge10^{20}$; this finite-sample fit is not an asymptotic tail
estimate. The largest recorded value among the $2\cdot10^7$ starts
is $\Phi_X/X=87.9$. Such samples see sets of
density down to about $10^{-7}$. A set that could spoil $T_X=o(X)$ needs
density only about $X^{-1}$, so the numerics support the model but do
not test the exceptional set.
\end{remark}
